\documentclass[10pt,a4paper]{amsart}
\usepackage[margin=19mm]{geometry}
\usepackage{amsmath,amssymb,mathtools}
\usepackage[scr=rsfs,cal=euler]{mathalfa}
\usepackage{xfrac}
\usepackage[unicode,hidelinks,bookmarks=false]{hyperref}
\newcommand{\ie}{i.e.\ }
\newcommand{\cf}{cf.\ }
\newcommand{\resp}{respectively\ }
\newcommand{\Subsection}{\subsection}
\providecommand{\nobreakdash}{\nobreak\hbox{-}}
\setlength{\emergencystretch}{2em}
\multlinegap0pt
\usepackage{amssymb}
\usepackage{mathtools}
\usepackage{xcolor}
\usepackage{tikz}
\usepackage[shortlabels]{enumitem}
\newtheorem{theorem}{Theorem}
\newtheorem{claim}{Claim}
\newcommand{\F} {\mathbb{F}}
\newcommand{\Height}    {\mathrm{h}}
\newcommand{\N} {\mathbb{N}}
\newcommand{\bigjoin}   {\bigvee}
\newcommand{\bigmeet}   {\bigwedge}
\title{Semirings of formal sums and injective partial transformations}
\author{Maximilien Gadouleau, Marianne Johnson}
\date{Source: arXiv:2603.26508v1 (CC BY 4.0)}
\begin{document}
\maketitle

\noindent\emph{Source and licence.} Semirings of formal sums and injective partial transformations. Version arXiv:2603.26508v1 (CC BY 4.0), distributed by arXiv under the Creative Commons Attribution 4.0 International licence, http://creativecommons.org/licenses/by/4.0/, as recorded in the arXiv licence metadata for this exact version and shown on the abstract page https://arxiv.org/abs/2603.26508v1. Full licence notice: \texttt{licenses/arxiv-2603.26508-CC-BY-4.0.txt}.\par\smallskip

\noindent\textit{Editorial scope.} Counted unit: the article's theorem theorem:division\_LC (source0.tex lines 1768--1770). The article's complete original proof of it is delivered with it (source0.tex lines 1773--1830, one \texttt{proof} environment of 58 lines). No mathematics is written, repaired or re-derived: the statement and the proof are the article's own lines, in the article's own notation, and no retained line is substituted or re-flowed.  Retained uncounted as the source-local context the proof needs: lines 1132--1146; lines 1725--1727.  Nothing was omitted from any retained span, so the package carries no omission record.  No retained line contains a citation, so the wrapper carries no bibliography and no bibliography entry.  No retained line contains a figure environment, an image-inclusion command or drawing code, so no image is delivered and the package declares no asset dependency.  Editorial formatting only, in addition: an AMS \texttt{amsart} preamble that restates the article's own declarations and macros used by the retained lines, namely the environments \texttt{theorem}, \texttt{claim} on separate counters, together with the standard packages those lines need.  The retained environments print in this document as Theorem 1, Claim 1 in order of appearance, whereas the article prints the same items with the numbering of its own full document.  The complete native source of the article as distributed in the arXiv e-print for this exact version is bundled unchanged as \texttt{sources/197289/source0.tex}.
\begin{theorem} \label{theorem:ax=b}
Let $a, b, x \in S$. For all $i \in \N$, let $l_i = b_i + a_0 b_i + a_i b_0$ and $u_i = l_i + a_i + 1$ and let
\begin{align*}
    \lambda_0 &= \bigjoin_{i \in \N} l_i \in S_0, \\
    \upsilon_0 &= \bigmeet_{i \in \N} u_i \in S_0.
\end{align*}
For all $i \ge 1$, let
\begin{align*}
    \lambda_i &= a_0 b_i + a_i b_0, \\
    \upsilon_i &= a_0 b_i + a_i b_0 + a_0 + 1.
\end{align*}

Then $ax = b$ if and only if $x_i \in [ \lambda_i, \upsilon_i ]$ for all $i \in \N$. 
Moreover, the set $\{ x \in S: ax = b \}$ is non-empty if and only if $\lambda_0 \le \upsilon_0$.
\end{theorem}

\begin{theorem} \label{theorem:division_Lh}
Let $a, b \in L^{(1)}$ with $\Height( a ) = h$. Then $x \in L^{(1)}$ satisfies $a x = b$ if and only if $x_{ -h } \in [b, b + a + L_h]$ in the Boolean algebra $V( \mathcal{L}_h )$.
\end{theorem}

\begin{theorem} \label{theorem:division_LC}
Let $a, b, x \in \F_2(L \cup C)$. Then, following the notation above, $ax = b$ if and only if $x^{(C)} \in X_t$ for some $t \in \{0,1\}$ and $x^{(\epsilon)} \in  Y_t^{(\epsilon)}$ for all $\epsilon \in \{0,1\}$.
\end{theorem}

\begin{proof}
We have $ax = b$ if and only if
\begin{align}
    \label{equation:bc}
    b^{(C)} &= a^{(C)} x^{(C)}, \\
    \label{equation:bepsilon}
    b^{(\epsilon)} &= a^{(\epsilon)} x^{(\epsilon)} + a^{(\epsilon)} x^{(C)} + a^{(C)} x^{(\epsilon)}, \forall \epsilon \in \{0,1\}.
\end{align}

Firstly, thanks to Theorem \ref{theorem:ax=b}, we can solve Equation \eqref{equation:bc}. 
As mentioned above, the set of $x^{(C)}$ solutions to Equation \eqref{equation:bc} with $|x^{(C)}| \equiv t \mod 2$ for each $t \in \{0,1\}$ is $X_t$.

Secondly, for $|x^{(C)}| \equiv t \mod 2$, Equation \eqref{equation:bepsilon} is equivalent to
\begin{equation}
    \label{equation:bepsilont}
    b^{(\epsilon)} = a^{(\epsilon)} x^{(\epsilon)} + t a^{(\epsilon)} + a^{(C)} x^{(\epsilon)}.
\end{equation}

\begin{claim}
Suppose $|a^{(C)}| \equiv 0 \mod 2$. Denoting $h = \Height( a^{(\epsilon)} )$, $x^{(\epsilon)}$ is a solution to Equation \eqref{equation:bepsilont} if and only if
\[
    x^{(\epsilon)}_{ -h } \in [ b^{(\epsilon)} + t a^{(\epsilon)}, b^{(\epsilon)} + (t + 1) a^{(\epsilon)} + L_h ]
\]
in $V( \mathcal{L}_h )$.
\end{claim}

\begin{proof}
If $|a^{(C)}| \equiv 0 \mod 2$, then Equation \eqref{equation:bepsilont} reduces to
\[
    a^{(\epsilon)} x^{(\epsilon)} = b^{(\epsilon)} + t a^{(\epsilon)},
\]
which by Theorem \ref{theorem:division_Lh}, yields $x^{(\epsilon)}_{ -h } \in [ b^{(\epsilon)} + t a^{(\epsilon)}, b^{(\epsilon)} + (t + 1) a^{(\epsilon)} + L_h ]$.
\end{proof}


\begin{claim}
Suppose $|a^{(C)}| \equiv 1 \mod 2$. Denoting $h = \max ( \Height( a^{(\epsilon)} ), \Height( b^{(\epsilon)} ) )$, $x^{(\epsilon)}$ is a solution to Equation \eqref{equation:bepsilont} if and only if
\[
    x^{(\epsilon)} = x^{(\epsilon)}_{ -h } \in [ b^{(\epsilon)} + t a^{(\epsilon)}, b^{(\epsilon)} + (t + 1) a^{(\epsilon)} ]
\]
in $V( \mathcal{L}_h )$.
\end{claim}


\begin{proof}
If $|a^{(C)}| \equiv 1 \mod 2$, then Equation \eqref{equation:bepsilont} reduces to
\[
    x^{(\epsilon)} = a^{(\epsilon)} x^{(\epsilon)} + b^{(\epsilon)} + t a^{(\epsilon)}.
\]
We note that this is not equivalent to $(a^{(\epsilon)} + 1)  x^{(\epsilon)} = b^{(\epsilon)} + t a^{(\epsilon)}$, for the multiplicative identity $1 = C_1$ does not belong to $L$. However, since the height of the right hand side is at most $h = \max ( \Height( a^{(\epsilon)} ), \Height( b^{(\epsilon)} ) )$, we have $\Height( x^{(\epsilon)} ) \le h$ and hence $x^{(\epsilon)} = x^{(\epsilon)}_{ -h }$.
The equation above is then equivalent to
\[
    ( a^{(\epsilon)} + L_h) x^{(\epsilon)} = b^{(\epsilon)} + t a^{(\epsilon)},
\]
which by Theorem \ref{theorem:division_Lh}, yields $x^{(\epsilon)} = x^{(\epsilon)}_{ -h } \in [ b^{(\epsilon)} + t a^{(\epsilon)}, b^{(\epsilon)} + (t + 1) a^{(\epsilon)} ]$.
\end{proof}

\end{proof}

\end{document}
